\documentclass{article}

\usepackage{booktabs}
\usepackage{tabularx}
\usepackage{hyperref}

\hypersetup{
    colorlinks=true,       % false: boxed links; true: colored links
    linkcolor=red,          % color of internal links (change box color with linkbordercolor)
    citecolor=green,        % color of links to bibliography
    filecolor=magenta,      % color of file links
    urlcolor=cyan           % color of external links
}

\title{Hazard Analysis\\\progname}

\author{\authname}

\date{}

\input{../Comments.text}
\input{../Common.text}

\begin{document}

\maketitle
\thispagestyle{empty}

~\newpage

\pagenumbering{roman}

\begin{table}[hp]
\caption{Revision History} \label{TblRevisionHistory}
\begin{tabularx}{\textwidth}{llX}
\toprule
\textbf{Date} & \textbf{Developer(s)} & \textbf{Change}\\
\midrule
Date1 & Name(s) & Description of changes\\
Date2 & Name(s) & Description of changes\\
... & ... & ...\\
\bottomrule
\end{tabularx}
\end{table}

~\newpage

\tableofcontents

~\newpage

\pagenumbering{arabic}


\section{Introduction}

MERLS is a telehealth platform for assessing the speech of bilingual children at home. 
Patients complete speech assessments online, and speech clinicians use the platform to review the recordings, verify the results, and manage their patients’ records and analytics. 
MERLS streamlines the clinician’s workflow by combining speech annotation and assessment analysis in one system. 
Annotation is supported by automatically generated transcripts that clinicians can review and correct, and analysis is supported by organized, visualized results that point out behaviors needing the clinician’s attention.

Because MERLS handles sensitive health information and supports clinical decisions about children, failures in the system could have serious consequences. 
A hazard is defined as any condition or state of the system that, under certain circumstances, could lead to harm to patients, such as an incorrect clinical conclusion or the disclosure of sensitive information. 
This document identifies the hazards in MERLS, analyzes how they could arise, and describes how the system’s design and requirements reduce them.


\section{Scope and Purpose of Hazard Analysis}

This document examines the hazards in each component of MERLS outlined below. 
For each hazard, the analysis determines the possible failures, along with their effects, causes, and recommended actions. 
It covers both technical hazards and user interaction hazards, because clinicians rely on the system’s outputs when assessing a child’s language ability, and the information it stores is confidential. 
If these hazards are not addressed, patients and clinicians could suffer losses such as the exposure of private health information, incorrect clinical conclusions that affect a child’s care, the permanent loss of assessment data, and a loss of confidence in the platform.

\section{System Boundaries and Components}

\wss{Dividing the system into components will help you brainstorm the hazards.
You shouldn't do a full design of the components, just get a feel for the major
ones.  For projects that involve hardware, the components will typically include
each individual piece of hardware.  If your software will have a database, or an
important library, these are also potential components.}

\section{Critical Assumptions}

\wss{These assumptions that are made about the software or system.  You should
minimize the number of assumptions that remove potential hazards.  For instance,
you could assume a part will never fail, but it is generally better to include
this potential failure mode.}

\section{Failure Mode and Effect Analysis}

\wss{Include your FMEA table here. This is the most important part of this document.}
\wss{The safety requirements in the table do not have to have the prefix SR.
The most important thing is to show traceability to your SRS. You might trace to
requirements you have already written, or you might need to add new
requirements.}
\wss{If no safety requirement can be devised, other mitigation strategies can be
entered in the table, including strategies involving providing additional
documentation, and/or test cases.}

\section{Safety and Security Requirements}

\wss{Newly discovered requirements.  These should also be added to the SRS.  (A
rationale design process how and why to fake it.)}

\section{Roadmap}

\wss{Which safety requirements will be implemented as part of the capstone timeline?
Which requirements will be implemented in the future?}

\newpage{}

\section{Appendix --- Generative AI Use Disclosure}


\subsection{AI Tools Used}

\begin{itemize}
    \item Claude Sonnet 5 (Medium)
    \item ChatGPT-5.6 Luna
\end{itemize}

\subsection{How and Where Used}

Both models were used in all the sections to do the following:
\begin{itemize}
    \item Check that each written section answered every question asked and met the rubric requirements
    \item Improve the wording, clarity and the flow of written sections
    \item Debug LaTeX compilation errors
\end{itemize}
The models were used to do the following in specific sections
\begin{itemize}
    \item Brainstorm possible failure modes for the system components in section 5
\end{itemize}


\subsection{Verification of AI-Generated Content}

To verify the output of the LLM, each team member went over their section and presented their work to the group, explaining what they had completed since the previous meeting. 
Each member was able to successfully explain their sections. 
Additionally, before any pull request was approved, a team member who had not written the section reviewed the work to ensure it was cohesive, accurate, and consistent with the rest of the document.

The factual content of the documents came from the research of the group members. 
The writer of each section researched, obtained, and documented all referenced sources. 
After the models suggested wording improvements, the writer re-read the section to confirm that it still conveyed the same point and that the research was not altered.


\newpage{}

\section{Appendix --- Reflection}


\input{../Reflection.text}

\begin{enumerate}
    \item What went well while writing this deliverable? 
    \item What pain points did you experience during this deliverable, and how
    did you resolve them?
    \item Which of your listed risks had your team thought of before this
    deliverable, and which did you think of while doing this deliverable? For
    the latter ones (ones you thought of while doing the Hazard Analysis), how
    did they come about?
    \item Other than the risk of physical harm (some projects may not have any
    appreciable risks of this form), list at least 2 other types of risk in
    software products. Why are they important to consider?
\end{enumerate}

\subsection*{Question 1}
\subsubsection*{Nour}
What went well with this deliverable was determining possible failure modes regarding the ML Transcription model. 
There are several things to be wary of when dealing with external APIs and using machine learning models that have been tuned for our case. 
This is because they are not always predictable. Models typically have some form of evaluation, and we were given the model already chosen by our supervisors and seeing that it has not been evaluated was a bit concerning since it is hard to be certain how well such a model will perform. 

What reassures me is that speech clinicians can go back and verify the model’s output against the video playback, which is extremely important. However, this places the responsibility on the clinician, including how much time they are willing to spend annotating and whether they will watch the entire assessment instead of skimming through it.

Another worrying issue is data retention. 
If Zero Data Retention was not applied for when the team started using OpenAI’s Whisper model, client information is being stored by OpenAI. 
OpenAI doesn’t use this data to train its models unless a user explicitly agrees to share it, but the data is still stored on its servers. 
OpenAI is a large and popular company that could experience cyber attacks, so having sensitive data stored in multiple places is not ideal.

\subsection*{Question 2}
This deliverable was relatively straightforward, and our team did not run into any major issues while completing it. 
We began by brainstorming the components of the system together, and we were able to identify what could go wrong with each one and possible hazards that could result. 
MERLS is a privacy-critical system, so security and data breaches were extremely important for us to address. 
The system also depends on performance and quality for different features to work, which made it easy to determine its high-level components. 
Because we share a strong understanding of the system, prioritizing the safety requirements was not difficult. 
Our supervisors also mentioned which parts of the system they wanted completed first, which gave us a clear direction as to what were the most important parts of the system.

\end{document}